\documentclass[../../main.tex]{subfiles}
\begin{document}
\section{丛同态}\label{section:ch10-丛同态}

若 $\pi\colon E\to M$ 与 $\pi'\colon E'\to M'$ 是向量丛，一个连续映射 $F\colon E\to E'$ 称为\textbf{丛同态}，若存在映射 $f\colon M\to M'$ 满足 $\pi'\circ F=f\circ\pi$，
\begin{align*}
\begin{array}{ccc}
E & \xrightarrow{F} & E'\\
\pi\downarrow && \downarrow\pi'\\
M & \xrightarrow{f} & M',
\end{array}
\end{align*}
且对每个 $p\in M$，限制映射 $F|_{E_p}\colon E_p\to E'_{f(p)}$ 是线性的。$F$ 与 $f$ 之间的关系表述为：$F$ \textbf{覆盖} $f$。

\begin{proposition}\label{proposition:ch10-丛同态覆盖映射}
设 $\pi\colon E\to M$ 与 $\pi'\colon E'\to M'$ 是向量丛，$F\colon E\to E'$ 是覆盖 $f\colon M\to M'$ 的丛同态。则 $f$ 连续且由 $F$ 唯一确定。若丛与 $F$ 都光滑，则 $f$ 也光滑。
\end{proposition}

\begin{proof}
所有结论都来自容易验证的事实 $f=\pi'\circ F\circ\zeta$，其中 $\zeta\colon M\to E$ 是零截面。
\end{proof}

一个双射的丛同态 $F\colon E\to E'$，若其逆也是丛同态，则称为\textbf{丛同构}；若 $F$ 还是微分同胚，则称为\textbf{光滑丛同构}。若 $E$ 与 $E'$ 之间存在（光滑）丛同构，则称这两个丛是（光滑）同构的。

在 $E$ 与 $E'$ 都是同一底空间 $M$ 上的向量丛这一特殊情形中，通常更有用的是一个稍强的丛同态概念。\textbf{$M$ 上的丛同态}是覆盖 $M$ 恒等映射的丛同态；换言之，是一个连续映射 $F\colon E\to E'$，满足 $\pi'\circ F=\pi$，
\begin{align*}
\begin{array}{ccc}
E & \xrightarrow{F} & E'\\
\pi\searrow && \swarrow\pi'\\
& M,
\end{array}
\end{align*}
且它在每条纤维上的限制是线性的。若存在 $M$ 上的丛同态 $F\colon E\to E'$ 同时又是（光滑）丛同构，则称 $E$ 与 $E'$ 在 $M$ 上（光滑）同构。下面的命题表明，不必检验逆映射的光滑性。

\begin{proposition}\label{proposition:ch10-双射丛同态}
设 $E$ 与 $E'$ 是带边或不带边光滑流形 $M$ 上的光滑向量丛，$F\colon E\to E'$ 是 $M$ 上的双射光滑丛同态。则 $F$ 是光滑丛同构。
\end{proposition}

\begin{proof}
习题~10-11。
\end{proof}

\begin{exercise}\label{exercise:ch10-27}
证明 $M$ 上的光滑 $k$ 秩向量丛是光滑平凡的，当且仅当它在 $M$ 上与乘积丛 $M\times\mathbb{R}^k$ 光滑同构。
\end{exercise}

\begin{example}[丛同态]\label{example:ch10-丛同态例子}
(a) 若 $F\colon M\to N$ 是光滑映射，则整体微分 $dF\colon TM\to TN$ 是覆盖 $F$ 的光滑丛同态。

(b) 若 $E\to M$ 是光滑向量丛，$S\subseteq M$ 是带边或不带边的浸入子流形，则包含映射 $E|_S\hookrightarrow E$ 是覆盖 $S$ 到 $M$ 包含映射的光滑丛同态。
\end{example}

设 $E\to M$ 与 $E'\to M$ 是带边或不带边光滑流形 $M$ 上的光滑向量丛，$\Gamma(E)$、$\Gamma(E')$ 为它们的光滑整体截面空间。若 $F\colon E\to E'$ 是 $M$ 上的光滑丛同态，则与 $F$ 复合诱导了一个映射 $\widetilde F\colon\Gamma(E)\to\Gamma(E')$，定义为
\begin{align}
\widetilde F(\sigma)(p)=(F\circ\sigma)(p)=F\big(\sigma(p)\big).\label{align:ch10-诱导截面映射}
\end{align}
容易验证 $\widetilde F(\sigma)$ 是 $E'$ 的截面，且由复合知其光滑。由于丛同态在纤维上是线性的，所得截面映射 $\widetilde F$ 在实数上是线性的。事实上，它满足更强的线性性质：若对任意光滑函数 $u_1,u_2\in C^\infty(M)$ 与光滑截面 $\sigma_1,\sigma_2\in\Gamma(E)$ 有
\begin{align*}
\widetilde F(u_1\sigma_1+u_2\sigma_2)=u_1\widetilde F(\sigma_1)+u_2\widetilde F(\sigma_2),
\end{align*}
则称映射 $\widetilde F\colon\Gamma(E)\to\Gamma(E')$ \textbf{在 $C^\infty(M)$ 上线性}。由定义~\eqref{align:ch10-诱导截面映射}~容易推知，光滑丛同态诱导的截面映射在 $C^\infty(M)$ 上线性。下面的引理表明逆命题也成立。

\begin{lemma}[丛同态刻画引理]\label{lemma:ch10-丛同态刻画}
设 $\pi\colon E\to M$ 与 $\pi'\colon E'\to M$ 为带边或不带边光滑流形 $M$ 上的光滑向量丛，$\Gamma(E)$、$\Gamma(E')$ 为其光滑截面空间。映射 $\widetilde F\colon\Gamma(E)\to\Gamma(E')$ 在 $C^\infty(M)$ 上线性，当且仅当存在 $M$ 上的光滑丛同态 $F\colon E\to E'$，使得对所有 $\sigma\in\Gamma(E)$ 有 $\widetilde F(\sigma)=F\circ\sigma$。
\end{lemma}

\begin{proof}
我们已指出，光滑丛同态诱导的截面映射在 $C^\infty(M)$ 上线性。反过来，假设 $\widetilde F\colon\Gamma(E)\to\Gamma(E')$ 在 $C^\infty(M)$ 上线性。首先证明 $\widetilde F$ 局部作用：若 $\sigma_1=\sigma_2$ 在某开集 $U\subseteq M$ 上成立，则 $\widetilde F(\sigma_1)=\widetilde F(\sigma_2)$ 在 $U$ 上成立。令 $\tau=\sigma_1-\sigma_2$；由 $\widetilde F$ 的线性性，只需假设 $\tau$ 在 $U$ 中为零并证明 $\widetilde F(\tau)$ 亦然。给定 $p\in U$，令 $\psi\in C^\infty(M)$ 为支撑在 $U$ 中、在 $p$ 处等于 $1$ 的光滑 bump 函数。由于 $\psi\tau$ 在 $M$ 上恒为零，$\widetilde F$ 在 $C^\infty(M)$ 上线性这一事实给出
\begin{align*}
0=\widetilde F(\psi\tau)=\psi\,\widetilde F(\tau).
\end{align*}
在 $p$ 处取值得 $\widetilde F(\tau)(p)=\psi(p)\widetilde F(\tau)(p)=0$；由于每个 $p\in U$ 都如此，结论成立。

接着证明 $\widetilde F$ 逐点作用：若 $\sigma_1(p)=\sigma_2(p)$，则 $\widetilde F(\sigma_1)(p)=\widetilde F(\sigma_2)(p)$。同样只需假设 $\tau(p)=0$ 并证 $\widetilde F(\tau)(p)=0$。令 $(\sigma_1,\dots,\sigma_k)$ 为 $E$ 在 $p$ 某邻域 $U$ 上的光滑局部标架，把 $\tau$ 按此标架写为 $\tau=u^i\sigma_i$，其中 $u^i$ 是 $U$ 上的光滑函数。$\tau(p)=0$ 意味着 $u^1(p)=\cdots=u^k(p)=0$。由向量丛与函数的扩张引理，存在光滑整体截面 $\widetilde\sigma_i\in\Gamma(E)$ 在 $p$ 附近与 $\sigma_i$ 一致，以及光滑函数 $\widetilde u^i\in C^\infty(M)$ 在 $p$ 附近与 $u^i$ 一致。于是在 $p$ 附近 $\tau=\widetilde u^i\widetilde\sigma_i$，故
\begin{align*}
\widetilde F(\tau)(p)=\widetilde F\big(\widetilde u^i\widetilde\sigma_i\big)(p)=\widetilde u^i(p)\widetilde F(\widetilde\sigma_i)(p)=0.
\end{align*}

如下定义丛同态 $F\colon E\to E'$：对任意 $p\in M$ 与 $v\in E_p$，令 $F(v)=\widetilde F(\widetilde v)(p)\in E'_p$，其中 $\widetilde v$ 是 $E$ 的任意光滑整体截面满足 $\widetilde v(p)=v$。上一段讨论表明所得 $E'_p$ 中元素与截面的选取无关。此映射显然满足 $\pi'\circ F=\pi$，且由 $\widetilde F$ 的线性性知它在每条纤维上线性。按定义它还满足 $F\circ\sigma(p)=\widetilde F(\sigma)(p)$ 对每个 $\sigma\in\Gamma(E)$。剩下只需证 $F$ 光滑；只需在每点附近证明。

给定 $p\in M$，令 $(\sigma_i)$ 为 $E$ 在 $p$ 某邻域上的光滑局部标架。由扩张引理，存在整体截面 $\widetilde\sigma_i$ 在 $p$ 的（更小的）邻域 $U$ 上与 $\sigma_i$ 一致。必要时再缩小 $U$，可设 $E'$ 在 $U$ 上有光滑局部标架 $(\sigma'_j)$。由于 $\widetilde F$ 把 $E$ 的光滑整体截面映为 $E'$ 的光滑整体截面，存在光滑函数 $A_i^j\in C^\infty(U)$ 使得 $\widetilde F(\widetilde\sigma_i)|_U=A_i^j\sigma'_j$。

对任意 $q\in U$ 与 $v\in E_q$，可写 $v=v^i\sigma_i(q)$（$v^1,\dots,v^k$ 为实数），于是
\begin{align*}
F\big(v^i\sigma_i(q)\big)=\widetilde F\big(v^i\widetilde\sigma_i\big)(q)=v^i\widetilde F(\widetilde\sigma_i)(q)=v^i A_i^j(q)\sigma'_j(q),
\end{align*}
因为 $v^i\widetilde\sigma_i$ 是 $E$ 在 $q$ 处取值为 $v$ 的光滑整体截面。若 $\Phi$ 与 $\Phi'$ 为分别与标架 $(\sigma_i)$、$(\sigma'_i)$ 相联系的 $E$、$E'$ 的局部平凡化，则复合映射 $\Phi'\circ F\circ\Phi^{-1}\colon U\times\mathbb{R}^k\to U\times\mathbb{R}^m$ 具有形式
\begin{align*}
\Phi'\circ F\circ\Phi^{-1}\big(q,(v^1,\dots,v^k)\big)=\big(q,(A_i^1(q)v^i,\dots,A_i^m(q)v^i)\big),
\end{align*}
它是光滑的。由于 $\Phi$ 与 $\Phi'$ 是微分同胚，这表明 $F$ 在 $\pi^{-1}(U)$ 上光滑。
\end{proof}

以后在有了更多工具后，我们会看到更多光滑丛同态的例子。眼下先给出几个初等例子。

\begin{example}[流形上的丛同态]\label{example:ch10-流形上丛同态}
(a) 若 $M$ 是光滑流形，$f\in C^\infty(M)$，则由 $X\mapsto fX$ 定义的 $\mathfrak X(M)$ 到自身的映射在 $C^\infty(M)$ 上线性，因为 $f(u_1X_1+u_2X_2)=u_1fX_1+u_2fX_2$，从而定义了 $M$ 上从 $TM$ 到自身的光滑丛同态。

(b) 若 $Z$ 是 $\mathbb{R}^3$ 上的光滑向量场，则与 $Z$ 的叉积定义了 $\mathfrak X(\mathbb{R}^3)$ 到自身的映射 $X\mapsto X\times Z$。由于它关于 $X$ 在 $C^\infty(\mathbb{R}^3)$ 上线性，它确定了 $\mathbb{R}^3$ 上从 $T\mathbb{R}^3$ 到 $T\mathbb{R}^3$ 的光滑丛同态。

(c) 给定 $Z\in\mathfrak X(\mathbb{R}^n)$，Euclid 点积定义了 $\mathfrak X(\mathbb{R}^n)$ 到 $C^\infty(\mathbb{R}^n)$ 的映射 $X\mapsto X\cdot Z$，它在 $C^\infty(\mathbb{R}^n)$ 上线性，从而确定了 $\mathbb{R}^n$ 上从 $T\mathbb{R}^n$ 到平凡线丛 $\mathbb{R}^n\times\mathbb{R}$ 的光滑丛同态。
\end{example}

由引理~\ref{lemma:ch10-丛同态刻画}，我们通常不再使用 $\widetilde F$ 这一记法，而用同一符号既表示 $M$ 上的丛同态 $F\colon E\to E'$，又表示它在截面上诱导的线性映射 $F\colon\Gamma(E)\to\Gamma(E')$；这两种类型的映射都称为丛同态。因为对截面的作用就是逐点施用丛同态，这不会引起混淆。事实上，我们在某些情形下一直都在这样做。例如，若 $a\in\mathbb{R}$，我们用同一记法 $X\mapsto aX$ 既表示在每个切空间 $T_pM$ 中用 $a$ 数乘向量，又表示用 $a$ 数乘向量场。由于数乘 $a$ 是从 $TM$ 到自身的丛同态，这不会有歧义。

应当注意，多数涉及微分的映射都不是丛同态。例如，若 $X$ 是光滑流形 $M$ 上的光滑向量场，Lie 导数算子 $\mathcal L_X\colon\mathfrak X(M)\to\mathfrak X(M)$ 不是从切丛到自身的丛同态，因为它不在 $C^\infty(M)$ 上线性。作为经验法则：把一个丛的光滑截面映到另一个丛的光滑截面的线性映射，若逐点作用则很可能是丛同态，但若涉及微分则不是。
\end{document}
